\documentclass{article}

% Ćirilica: kodiranje fonta i jezik (srpski, ćirilica)
\usepackage[T2A]{fontenc}
\usepackage[utf8]{inputenc}
\usepackage[serbianc]{babel}

% Veličina strane i margine
\usepackage[a4paper,top=2cm,bottom=2cm,left=3cm,right=3cm,marginparwidth=1.75cm]{geometry}

% Korisni paketi
\usepackage{amsmath}
\usepackage{graphicx}
\usepackage{float}
\usepackage{tcolorbox}
\usepackage[colorlinks=true, allcolors=blue]{hyperref}

% Цртежи (TikZ)
\usepackage{tikz}
\usetikzlibrary{arrows.meta, calc, decorations.markings, babel}

% Електричне шеме (отпорници, кондензатори, извори)
\usepackage[european]{circuitikz}

% TikZ стилови који се понављају на цртежима
\tikzset{
    naboj/.style={circle, draw, thick, minimum size=6mm, inner sep=0pt, font=\small},
    jezgro/.style={circle, draw, thick, fill=gray!15, minimum size=7mm, inner sep=0pt, font=\small},
    elektron/.style={circle, draw, fill=cyan!30, inner sep=0pt, minimum size=5pt},
    supljina/.style={circle, draw, thick, minimum size=5mm, inner sep=0pt, font=\small},
    sredstrelica/.style={postaction={decorate}, decoration={markings, mark=at position 0.55 with {\arrow{Stealth}}}},
}

\title{ОЕТ 1 - Струја и колу}
\author{Саобраћајно техничка школа, Београд, Земун}

\begin{document}
\maketitle

\section{Компоненте у колу}

% ------------------------------------------------------------------
\subsection{Једносмерна и наизменична струја}

Струја је усмерено кретање наелектрисања. Према томе како се мења током времена, делимо је на \textbf{једносмерну} и \textbf{наизменичну}.

\textbf{Једносмерна струја (DC)} тече увек у истом смеру. Најчешће јој је и јачина стална. Добијамо је из батерије, акумулатора или исправљача (пуњача).

\textbf{Наизменична струја (AC)} периодично мења и смер и јачину. Најчешће има облик синусоиде. Таква је струја у градској мрежи.

\begin{figure}[H]
    \centering
    \begin{tikzpicture}[>=Stealth]
        % --- (а) једносмерна ---
        \draw[->] (0,0) -- (4.2,0) node[below] {$t$};
        \draw[->] (0,-1.4) -- (0,1.6) node[left] {$i$};
        \draw[very thick, blue] (0,0.9) -- (3.9,0.9);
        \node[left] at (0,0.9) {$I$};
        \node at (2,-1.9) {(а) једносмерна};
        % --- (б) наизменична ---
        \begin{scope}[xshift=5.6cm]
            \draw[->] (0,0) -- (4.2,0) node[below] {$t$};
            \draw[->] (0,-1.4) -- (0,1.6) node[left] {$i$};
            \draw[very thick, red, domain=0:3.9, samples=120] plot (\x, {sin(180*\x/1.5)});
            % период T
            \draw[dashed] (0.375,1) -- (0.375,1.25);
            \draw[dashed] (1.875,1) -- (1.875,1.25);
            \draw[<->] (0.375,1.2) -- (1.875,1.2) node[midway, above] {$T$};
            \node at (2,-1.9) {(б) наизменична};
        \end{scope}
    \end{tikzpicture}
    \caption{Струја током времена: (а) стално у истом смеру, (б) мења смер}
    \label{fig:dc-ac}
\end{figure}

Наизменичну струју описујемо са две величине:
\begin{itemize}
    \item \textbf{Период} $T$ -- време за које се заврши један цео циклус. Јединица је секунда (s).
    \item \textbf{Фреквенција} $f$ -- колико се циклуса деси у једној секунди. Јединица је херц (Hz).
\end{itemize}

\begin{equation}
    f = \frac{1}{T}
\end{equation}

У градској мрежи напон је 230~V, а фреквенција $f = 50$~Hz. Значи да је $T = \frac{1}{50}~\mathrm{s} = 20$~ms и да струја промени смер 100 пута у секунди.

Код наизменичне струје електрони не путују дуж жице. Они само осцилују напред-назад, а енергија се ипак преноси до потрошача.

\begin{figure}[H]
    \centering
    \begin{circuitikz}
        \draw (0,2) to[battery1] (0,0);
        \node[right] at (0.35,1.35) {$+$};
        \node[right] at (0.35,0.65) {$-$};
        \node at (0,-0.6) {једносмерни извор};
        \draw (4,0) to[sV] (4,2);
        \node at (4,-0.6) {наизменични извор};
    \end{circuitikz}
    \caption{Ознаке извора у шемама}
    \label{fig:izvori}
\end{figure}

% ------------------------------------------------------------------
\subsection{Отпорност и Омов закон}

Сваки материјал се мање или више опире протицању струје. Електрони се при кретању сударају са атомима материјала и то их успорава. Та особина се зове \textbf{електрична отпорност} и обележава се са $R$. Због тих судара се материјал и загрева.

Елементи код којих је отпорност изражена зову се \textbf{отпорници}. На шемама се обележавају на један од следећих начина:

\begin{figure}[H]
    \centering
    \begin{circuitikz}
        \draw (0,0) to[american resistor, l=$R$, o-o] (3,0);
        % меандар
        \begin{scope}[xshift=4cm]
            \draw (0,0) circle (1.5pt);
            \draw (0,0) -- (0.6,0) -- (0.6,0.3) -- (0.8,0.3) -- (0.8,-0.3) -- (1.0,-0.3) -- (1.0,0.3)
                -- (1.2,0.3) -- (1.2,-0.3) -- (1.4,-0.3) -- (1.4,0.3) -- (1.6,0.3) -- (1.6,-0.3)
                -- (1.8,-0.3) -- (1.8,0.3) -- (2.0,0.3) -- (2.0,-0.3) -- (2.2,-0.3) -- (2.2,0) -- (3,0);
            \draw (3,0) circle (1.5pt);
            \node at (1.5,0.65) {$R$};
        \end{scope}
        \draw (8,0) to[european resistor, l=$R$, o-o] (11,0);
    \end{circuitikz}
    \caption{Симболи отпорника}
    \label{fig:simboli-otpornika}
\end{figure}

\subsubsection{Од чега зависи отпорност}

Отпорност жице зависи од три ствари:
\begin{itemize}
    \item \textbf{дужине} $l$: дужа жица има већу отпорност, јер електрони прелазе дужи пут;
    \item \textbf{површине попречног пресека} $S$: дебља жица има мању отпорност, јер електрони имају више места;
    \item \textbf{материјала}: то описује \textbf{специфична отпорност} $\rho$.
\end{itemize}

\begin{equation}
    R = \rho \frac{l}{S}
\end{equation}

\begin{figure}[H]
    \centering
    \begin{tikzpicture}[>=Stealth]
        % жица (ваљак)
        \draw[thick] (0,0.5) -- (6,0.5);
        \draw[thick] (0,-0.5) -- (6,-0.5);
        \draw[thick, fill=gray!20] (0,0) ellipse (0.25 and 0.5);
        \draw[thick] (6,0.5) arc (90:-90:0.25 and 0.5);
        \draw[thick, dashed] (6,0.5) arc (90:270:0.25 and 0.5);
        \node at (0,0) {$S$};
        % дужина
        \draw[<->] (0,-0.9) -- (6,-0.9) node[midway, below] {$l$};
    \end{tikzpicture}
    \caption{Проводник дужине $l$ и попречног пресека $S$}
    \label{fig:zica}
\end{figure}

Јединица специфичне отпорности је $\rho_u = \Omega \cdot \mathrm{m}$. Мање $\rho$ значи бољи проводник:

\begin{table}[H]
    \centering
    \begin{tabular}{|l|c|}
        \hline
        Материјал & $\rho$ на $20^\circ$C [$\Omega \cdot \mathrm{m}$] \\
        \hline
        сребро & $1{,}59 \cdot 10^{-8}$ \\
        бакар & $1{,}68 \cdot 10^{-8}$ \\
        алуминијум & $2{,}65 \cdot 10^{-8}$ \\
        гвожђе & $9{,}7 \cdot 10^{-8}$ \\
        константан & $4{,}9 \cdot 10^{-7}$ \\
        \hline
    \end{tabular}
    \caption{Специфична отпорност неких материјала}
    \label{tab:rho}
\end{table}

Каблови се праве од бакра: проводи скоро као сребро, а много је јефтинији.

\subsubsection{Омов закон}

Да би кроз проводник текла стална струја, на његовим крајевима мора да постоји сталан напон. Ако се напон промени, сразмерно се мења и струја. Ту везу је експериментално утврдио немачки физичар Ом, па се зове \textbf{Омов закон}:

\begin{equation}
    I = \frac{U}{R}
\end{equation}

Из њега следи и:

\begin{equation}
    U = R \cdot I \qquad \text{и} \qquad R = \frac{U}{I}
\end{equation}

\begin{figure}[H]
    \centering
    \begin{circuitikz}[>=Stealth]
        \draw (0,3) to[battery1] (0,0);
        \draw (0,3) -- (1,3)
              to[european resistor, l=$R$] (4,3) -- (4,0) -- (0,0);
        \node[left] at (-0.35,1.8) {$+$};
        \node[left] at (-0.35,1.2) {$-$};
        % струја: из + извора, кроз отпорник с лева на десно
        \draw[->, thick] (1.5,2.4) -- (3.5,2.4) node[midway, below] {$I$};
        % ознака напона на отпорнику: показује на +
        \draw[->, thick] (3.6,3.6) to[out=120, in=60] (1.4,3.6);
        \node at (2.5,4.45) {$U$};
        \node at (1.1,3.55) {$+$};
    \end{circuitikz}
    \caption{Најједноставније коло: извор и отпорник}
    \label{fig:om-kolo}
\end{figure}

Јединица отпорности је \textbf{ом}:

\begin{equation}
    R_u = \frac{U_u}{I_u} = \frac{\mathrm{V}}{\mathrm{A}} = \text{ом} = \Omega
\end{equation}

Отпорност од једног ома има проводник кроз који при напону од 1~V тече струја од 1~A.

\begin{tcolorbox}[colback=blue!5, colframe=blue!60, title=Омов троугао]
    \begin{minipage}{0.3\linewidth}
        \centering
        \begin{tikzpicture}
            \draw[thick] (0,0) -- (2.4,0) -- (1.2,2) -- cycle;
            \draw[thick] (0.6,1) -- (1.8,1);
            \draw[thick] (1.2,0) -- (1.2,1);
            \node at (1.2,1.35) {$U$};
            \node at (0.75,0.45) {$I$};
            \node at (1.65,0.45) {$R$};
        \end{tikzpicture}
    \end{minipage}
    \hfill
    \begin{minipage}{0.65\linewidth}
        Покриј прстом величину коју тражиш, а оно што остане је формула:
        \begin{itemize}
            \item покријеш $U$: остаје $I \cdot R$
            \item покријеш $I$: остаје $\frac{U}{R}$
            \item покријеш $R$: остаје $\frac{U}{I}$
        \end{itemize}
    \end{minipage}
\end{tcolorbox}

\textbf{Пример:} батерија од $U = 12$~V прикључена је на отпорник $R = 4~\Omega$. Струја је $I = \frac{U}{R} = \frac{12~\mathrm{V}}{4~\Omega} = 3$~A.

% ------------------------------------------------------------------
\subsection{Капацитивност и кондензатор}

\textbf{Кондензатор} чине два проводника (\textbf{електроде}) раздвојена изолатором (\textbf{диелектриком}). Према облику електрода може бити плочасти, цилиндрични или сферни. Према диелектрику може бити ваздушни, керамички, папирни\dots

Кад на кондензатор доведемо напон $U$, на једној електроди се скупи наелектрисање $+Q$, а на другој $-Q$. Колико наелектрисања кондензатор прими при неком напону, показује \textbf{капацитивност} $C$:

\begin{equation}
    C = \frac{Q}{U}
\end{equation}

Јединица капацитивности је \textbf{фарад}:

\begin{equation}
    C_u = \frac{Q_u}{U_u} = \frac{\mathrm{C}}{\mathrm{V}} = \text{фарад} = \mathrm{F}
\end{equation}

Један фарад је врло велика јединица. У пракси се користе мање: микрофарад ($1~\mu\mathrm{F} = 10^{-6}$~F), нанофарад ($1~\mathrm{nF} = 10^{-9}$~F) и пикофарад ($1~\mathrm{pF} = 10^{-12}$~F).

\begin{figure}[H]
    \centering
    \begin{circuitikz}
        \draw (0,0) to[C, l=$C$, o-o] (2.5,0);
        \node at (1.25,-1) {(а) стална};
        \draw (4,0) to[vC, l=$C$, o-o] (6.5,0);
        \node at (5.25,-1) {(б) променљива};
    \end{circuitikz}
    \caption{Ознаке кондензатора: (а) сталне и (б) променљиве капацитивности}
    \label{fig:simboli-kondenzatora}
\end{figure}

\subsubsection{Плочасти кондензатор}

\begin{figure}[H]
    \centering
    \begin{tikzpicture}[>=Stealth]
        % задња плоча (+Q)
        \draw[thick, fill=red!8] (0,0) -- (0,2.6) -- (1.1,3.4) -- (1.1,0.8) -- cycle;
        % предња плоча (-Q)
        \draw[thick, fill=blue!8] (1,0) -- (1,2.6) -- (2.1,3.4) -- (2.1,0.8) -- cycle;
        \node at (0.35,3.35) {$+Q$};
        \node at (2.45,3.55) {$-Q$};
        \node at (1.65,1.5) {$S$};
        % растојање d
        \draw[<->] (0,-0.3) -- (1,-0.3) node[midway, below] {$d$};
    \end{tikzpicture}
    \caption{Плочасти кондензатор: површина плоча $S$, растојање $d$}
    \label{fig:plocasti}
\end{figure}

Капацитивност плочастог кондензатора у ваздуху (вакууму) је:

\begin{equation}
    C_0 = \varepsilon_0 \frac{S}{d}
\end{equation}

Ако се између плоча стави диелектрик, капацитивност порасте онолико пута колико је $\varepsilon_r$:

\begin{equation}
    C = \varepsilon_r C_0 = \varepsilon \frac{S}{d}
\end{equation}

Капацитивност је дакле већа ако су:
\begin{itemize}
    \item плоче \textbf{веће} (веће $S$),
    \item плоче \textbf{ближе} (мање $d$),
    \item између њих диелектрик са \textbf{већим} $\varepsilon_r$ (видети табелу \ref{tab:dielektrici}).
\end{itemize}

Зато се кондензатори праве од танких металних фолија које су веома близу, раздвојене чврстим диелектриком, и смотане у ролну.

Као кондензатор се понашају и сви метални делови раздвојени изолацијом, ако између њих постоји напон. На пример, две жице у истом каблу.

\subsubsection{Кондензатор у колу једносмерне струје}

\begin{figure}[H]
    \centering
    \begin{circuitikz}[>=Stealth]
        \draw (0,3) to[battery1] (0,0);
        \draw (0,3) -- (1,3)
              to[european resistor, l=$R$] (3.5,3) -- (4.5,3)
              to[C, l=$C$] (4.5,0) -- (0,0);
        \node[left] at (-0.35,1.8) {$+$};
        \node[left] at (-0.35,1.2) {$-$};
        \draw[->, thick] (1.2,2.4) -- (3.3,2.4) node[midway, below] {$I$};
    \end{circuitikz}
    \caption{Кондензатор у колу једносмерне струје}
    \label{fig:c-dc-kolo}
\end{figure}

\begin{itemize}
    \item У тренутку укључења кондензатор је празан и струја тече.
    \item Електрони долазе на једну плочу, а одлазе са друге, привидно је као да \textbf{прескачу} на другу страну. Кроз диелектрик не пролазе.
    \item Како се кондензатор пуни, напон на њему расте, а струја опада.
\end{itemize}

\begin{figure}[H]
    \centering
    \begin{tikzpicture}[>=Stealth]
        \draw[->] (0,0) -- (7,0) node[below] {$t$};
        \draw[->] (0,0) -- (0,3) ;
        % напон извора
        \draw[dashed] (0,2.4) -- (6.6,2.4);
        \node[left] at (0,2.4) {$U$};
        % напон на кондензатору
        \draw[very thick, blue, domain=0:6.5, samples=100] plot (\x, {2.4*(1-exp(-\x/1.2))});
        \node[blue, right] at (6.5,2.25) {$u_C$};
        % струја
        \draw[very thick, red, domain=0:6.5, samples=100] plot (\x, {2.1*exp(-\x/1.2)});
        \node[red, right] at (6.5,0.25) {$i$};
        % тау и 5 тау
        \draw[dotted] (1.2,0) -- (1.2,1.52);
        \node[below] at (1.2,0) {$\tau$};
        \draw[dotted] (6,0) -- (6,2.4);
        \node[below] at (6,0) {$5\tau$};
    \end{tikzpicture}
    \caption{Пуњење кондензатора: напон $u_C$ расте, струја $i$ опада}
    \label{fig:punjenje}
\end{figure}

Напуњен кондензатор задржава наелектрисање. Ако га откачимо од извора и вежемо на неки потрошач, он се празни кроз њега, као мала батерија. Тако ради блиц на фотоапарату.

\begin{tcolorbox}[colback=blue!5, colframe=blue!60, title=Важно]
    У колу једносмерне струје напуњен кондензатор \textbf{не пропушта струју}. Понаша се као \textbf{прекид}.
\end{tcolorbox}

\subsubsection{Кондензатор у колу наизменичне струје}

\begin{itemize}
    \item Напон извора стално мења смер, па се кондензатор стално пуни, празни и пуни у супротном смеру.
    \item Зато у колу \textbf{стално тече струја}, иако кроз диелектрик не прође ни један електрон.
    \item Ако редно са кондензатором вежемо сијалицу, она ће стално светлети. На батерији би само кратко бљеснула.
\end{itemize}

\begin{figure}[H]
    \centering
    \begin{circuitikz}
        \draw (0,0) to[sV] (0,3) -- (1,3)
              to[C, l=$C$] (3.5,3) -- (4.5,3)
              to[lamp] (4.5,0) -- (0,0);
    \end{circuitikz}
    \caption{Сијалица редно са кондензатором у колу наизменичне струје светли}
    \label{fig:c-ac-kolo}
\end{figure}

\begin{tcolorbox}[colback=blue!5, colframe=blue!60, title=Важно]
    \begin{itemize}
        \item \textbf{Једносмерна струја}: кондензатор се напуни и после тога \textbf{не пропушта} струју (прекид).
        \item \textbf{Наизменична струја}: кондензатор се стално пуни и празни, па \textbf{пропушта} струју.
        \item Што је фреквенција \textbf{већа}, кондензатор \textbf{боље} пропушта струју.
    \end{itemize}
\end{tcolorbox}

\end{document}